\chapter{MARCO METODOLÓGICO}

Este capítulo describe la metodología experimental de DARL. El objetivo metodológico
es construir episodios controlados de drift sobre datos tabulares reales,
observar el efecto del cambio mediante métricas estadísticas y supervisadas, y
entrenar una política que seleccione la acción de mantenimiento más eficiente
para un \textit{pipeline} de dos etapas.

% ------------------------------------------------------------
\section{Diseño general del framework DARL}
% ------------------------------------------------------------

DARL se organiza como un ciclo de seis módulos. Primero, se seleccionan
\textit{datasets} tabulares de TableShift y se entrena un \textit{pipeline} base
con XGBoost como modelo principal. Segundo, se aplica drift sintético controlado
sobre variables numéricas, categóricas o sobre las etiquetas. Tercero, se
monitorea el cambio resultante mediante señales de \textit{covariate shift},
\textit{concept drift} y degradación de desempeño. Cuarto, un agente de
aprendizaje por refuerzo observa esas señales y elige una acción de
mantenimiento. Quinto, se ejecuta la acción seleccionada y se mide su efecto en
rendimiento y cómputo. Finalmente, esas mediciones alimentan la recompensa que
guía el entrenamiento de PPO.

\begin{figure}[H]
    \centering
    \resizebox{0.98\textwidth}{!}{\begin{tikzpicture}[
    font=\small,
    stagebox/.style={
        draw,
        rounded corners=2pt,
        fill=gray!8,
        minimum height=24mm,
        text width=27mm,
        align=center
    },
    arrow/.style={-{Latex[length=2.3mm]}, thick},
    feedback/.style={-{Latex[length=2.3mm]}, thick, dashed}
]

\node[stagebox] (s1) at (0, 0) {\textbf{1. Datos}\\TableShift\\XGBoost};
\node[stagebox, right=6mm of s1] (s2) {\textbf{2. Drift}\\Beta numérica\\remuestreo categórico\\label flipping};
\node[stagebox, right=6mm of s2] (s3) {\textbf{3. Monitoreo}\\Covariate shift\\Concept drift};
\node[stagebox, right=6mm of s3] (s4) {\textbf{4. Acción RL}\\Defer\\Update features\\Update model\\Retrain all};
\node[stagebox, right=6mm of s4] (s5) {\textbf{5. Métricas}\\$\Delta$ AUC-ROC\\F1\\tiempo y RAM};
\node[stagebox, right=6mm of s5] (s6) {\textbf{6. PPO}\\reward\\o punishment};

\draw[arrow] (s1) -- (s2);
\draw[arrow] (s2) -- (s3);
\draw[arrow] (s3) -- (s4);
\draw[arrow] (s4) -- (s5);
\draw[arrow] (s5) -- (s6);
\draw[feedback] (s6.south) -- ++(0,-9mm) -| (s4.south);
\draw[feedback] (s6.south) -- ++(0,-16mm) -| (s2.south);

\end{tikzpicture}
}
    \caption{Flujo metodológico propuesto para DARL}
    \label{fig:metodologia_darl}
\end{figure}

Como se muestra en la Figura~\ref{fig:metodologia_darl}, la metodología no trata
la detección de drift como un fin aislado. Las métricas de monitoreo son señales
intermedias para decidir si conviene diferir, actualizar el preprocesamiento,
actualizar el modelo predictivo o reentrenar todo el \textit{pipeline}. Esta
separación responde a la brecha identificada en la literatura: diagnosticar que
existe drift determina qué componente debe actualizarse
\cite{cai2023disde,mahadevan2024cost}.

% ------------------------------------------------------------
\section{Selección de datasets y particiones experimentales}
% ------------------------------------------------------------

La evaluación se concentra en dos \textit{datasets} de TableShift
\cite{gardner2023tableshift}. La selección busca cubrir dos regímenes tabulares
distintos: uno dominado por variables categóricas y otro dominado por variables
numéricas. Esta decisión permite evaluar si las acciones de actualización
selectiva se comportan de forma distinta cuando el drift afecta principalmente
frecuencias categóricas o distribuciones numéricas continuas.

\begin{table}[H]
    \centering
    \caption{Datasets considerados en la metodología DARL}
    \label{tab:datasets_metodologia_darl}
    \resizebox{0.98\textwidth}{!}{%
        \begin{tabular}{lll}
            \hline
            \textbf{Dataset}             & \textbf{Rol experimental} & \textbf{Tipo dominante de variables} \\
            \hline
            \text{diabetes\_readmission} & Readmisión hospitalaria   & Categóricas y discretas              \\
            \text{physionet}             & Predicción de sepsis      & Numéricas fisiológicas               \\
            \hline
        \end{tabular}
    }
\end{table}

Para cada dataset se distinguen tres particiones. La partición de referencia se
usa para ajustar el preprocesamiento y entrenar el modelo base. La partición
pre-drift estima el desempeño inicial sin intervención. La partición post-drift
recibe las perturbaciones sintéticas y permite medir degradación, detección y
recuperación. En el caso de \text{diabetes\_readmission}, el objetivo binario
representa si un encuentro hospitalario termina en readmisión frente a la clase
``NO''. En el caso de \text{physionet}, el objetivo binario representa la
etiqueta de sepsis usada en los notebooks exploratorios del repositorio.

Esta estructura conserva dos propiedades metodológicas importantes. Primero,
usa datos tabulares reales y particiones reproducibles, en lugar de construir
todo el entorno con datos sintéticos. Segundo, mantiene control experimental
sobre el tipo y severidad de drift, porque el cambio se inyecta después de fijar
la distribución de referencia.

% ------------------------------------------------------------
\section{Pipeline base de dos etapas}
% ------------------------------------------------------------

El sistema evaluado es un \textit{pipeline} de dos etapas. La primera etapa transforma los datos crudos en
una representación numérica para el modelo. La segunda etapa aprende el mapeo
predictivo. La implementación metodológica considera XGBoost como modelo
principal por su desempeño en datos tabulares \cite{chen2016xgboost}, y
regresión logística como línea base de menor complejidad.

La etapa de preprocesamiento aplica tres familias de transformación. Las
variables numéricas sesgadas se transforman mediante cuantiles, las variables
categóricas de alta cardinalidad se codifican mediante \textit{target encoding},
y el resto de variables numéricas se escala con una normalización estándar. La
separación entre preprocesamiento y modelo es deliberada: permite evaluar
acciones que actualizan solo los parámetros de transformación, solo el modelo
predictivo o ambos.

En cada configuración experimental, el \textit{pipeline} base se ajusta con la
partición de referencia. Luego se evalúa en la partición pre-drift para obtener
el desempeño inicial y en la partición post-drift perturbada para estimar la
degradación. El valor de referencia para la recompensa no es el rendimiento
absoluto del modelo, sino la recuperación lograda por cada acción frente al
estado drifted sin mantenimiento.

% ------------------------------------------------------------
\section{Inyección sintética de drift}
% ------------------------------------------------------------

El módulo \text{DriftInjector} genera tres escenarios: \textit{covariate
    shift}, \textit{concept drift} y la combinación de ambos. El drift se controla
mediante una severidad $\gamma \in [0,1]$, donde valores mayores aumentan la
proporción de observaciones afectadas o la intensidad del desplazamiento. Esta
formulación mantiene trazabilidad entre el tipo de perturbación aplicada y las
señales que deberían observarse en el módulo de monitoreo.

\subsection{Drift numérico mediante mezcla Beta}

Para variables numéricas, primero se normaliza cada valor al intervalo $[0,1]$
usando los mínimos y máximos de referencia. La distribución objetivo que genera
los valores perturbados se define mediante una Beta:

\begin{equation}
    z_{ij} \sim \mathrm{Beta}(\alpha_j,\beta_j),
    \qquad
    f_j(z)=
    \frac{z^{\alpha_j-1}(1-z)^{\beta_j-1}}
    {B(\alpha_j,\beta_j)},
    \quad 0 \leq z \leq 1.
    \label{eq:beta_target_distribution}
\end{equation}

\noindent\textbf{Donde:}
\begin{itemize}
    \item $z_{ij}$ representa el valor sintético normalizado generado para la observación $i$ en la variable numérica $j$.
    \item $\mathrm{Beta}(\alpha_j,\beta_j)$ representa la distribución objetivo usada para generar valores en $[0,1]$.
    \item $f_j(z)$ representa la función de densidad de la distribución Beta asociada a la variable $j$.
    \item $\alpha_j$ y $\beta_j$ son parámetros positivos que controlan la forma de la distribución objetivo.
    \item $B(\alpha_j,\beta_j)$ representa la función Beta de normalización.
\end{itemize}

Conceptualmente, la distribución Beta no decide por sí sola qué filas son
modificadas; solo define de dónde provienen los valores numéricos alternativos.
Al variar $\alpha_j$ y $\beta_j$, el inyector puede generar desplazamientos
hacia valores altos, bajos, centrales o extremos dentro del rango normalizado.

La selección de las observaciones perturbadas se modela por separado mediante
una variable Bernoulli:

\begin{equation}
    m_{ij} \sim \mathrm{Bernoulli}(\gamma),
    \qquad
    \tilde{x}_{ij}
    =
    \left[
        (1-m_{ij})u_{ij}
        + m_{ij}z_{ij}
        \right]
    (x^{\max}_{j}-x^{\min}_{j})
    + x^{\min}_{j}.
    \label{eq:beta_mixture_drift}
\end{equation}

\noindent\textbf{Donde:}
\begin{itemize}
    \item $m_{ij}$ representa la máscara binaria que indica si la observación $i$ es perturbada en la variable $j$.
    \item $\gamma$ representa la severidad del drift y controla la probabilidad de reemplazo.
    \item $u_{ij}$ representa el valor original normalizado al intervalo $[0,1]$.
    \item $z_{ij}$ representa la muestra Beta definida en la Ecuación~\ref{eq:beta_target_distribution}.
    \item $\tilde{x}_{ij}$ representa el valor de la variable numérica $j$ para la observación $i$ después de aplicar drift.
    \item $x^{\min}_{j}$ y $x^{\max}_{j}$ son el mínimo y máximo de la variable $j$ en la partición de referencia.
\end{itemize}

Conceptualmente, la Bernoulli controla la intensidad aplicada a nivel de
observación y variable: si $m_{ij}=0$, se conserva el valor original; si
$m_{ij}=1$, se usa el valor sintético generado por la Beta. Así, el mecanismo
simula \textit{covariate shift} en $P(X)$ sin cambiar directamente la etiqueta,
por lo que debería reflejarse principalmente en métricas como KS, PSI o C2ST.

\subsection{Drift categórico mediante remuestreo marginal}

Para variables categóricas, se construye una distribución objetivo $q_j(c)$ sobre
las categorías de la variable $j$ y se interpola con la distribución de
referencia $p_{0,j}(c)$:

\begin{equation}
    p_{\gamma,j}(c)
    =
    (1-\gamma)p_{0,j}(c)
    + \gamma q_j(c).
    \label{eq:categorical_marginal_drift}
\end{equation}

\noindent\textbf{Donde:}
\begin{itemize}
    \item $p_{\gamma,j}(c)$ representa la probabilidad post-drift de la categoría $c$ en la variable categórica $j$.
    \item $p_{0,j}(c)$ representa la frecuencia relativa de la categoría $c$ en la partición de referencia.
    \item $q_j(c)$ representa la distribución categórica objetivo definida para inducir drift.
    \item $\gamma$ representa la severidad del drift y controla cuánto se acerca la distribución post-drift a $q_j(c)$.
    \item $c$ representa una categoría posible de la variable $j$.
\end{itemize}

Conceptualmente, esta fórmula modifica las frecuencias de las categorías sin
alterar la regla de etiquetado. En la implementación, esta distribución se
aplica mediante remuestreo de filas ponderado por las nuevas probabilidades
marginales. Este diseño es más conservador que reemplazar categorías de manera
independiente, porque reduce el riesgo de crear combinaciones de columnas poco
plausibles.

\subsection{Concept drift mediante label flipping}

Para simular \textit{concept drift}, se altera la relación entre $X$ e $Y$
mediante inversión aleatoria de etiquetas binarias. La probabilidad de inversión
se acota para evitar un escenario completamente aleatorio:

\begin{equation}
    \tilde{y}_{i}
    =
    \left\{
    \begin{array}{ll}
        1-y_i, & \text{con probabilidad } p_{\text{flip}}=\min(0.45,0.45\gamma), \\
        y_i,   & \text{con probabilidad } 1-p_{\text{flip}}.
    \end{array}
    \right.
    \label{eq:label_flipping}
\end{equation}

\noindent\textbf{Donde:}
\begin{itemize}
    \item $\tilde{y}_{i}$ representa la etiqueta de la observación $i$ después de aplicar drift.
    \item $y_i$ representa la etiqueta original antes de aplicar drift.
    \item $p_{\text{flip}}$ representa la probabilidad de invertir la etiqueta.
    \item $\gamma$ representa la severidad del drift.
    \item $0.45$ es el límite máximo de inversión usado para mantener una señal predictiva remanente.
\end{itemize}

Conceptualmente, el \textit{label flipping} no cambia la distribución marginal de
las covariables $P(X)$, pero sí degrada la relación condicional $P(Y \mid X)$.
Por ello, este mecanismo permite evaluar si el sistema distingue escenarios en
los que actualizar el preprocesamiento no basta y se requiere actualizar el
modelo predictivo.

% ------------------------------------------------------------
\section{Monitoreo y señales de diagnóstico}
% ------------------------------------------------------------

El módulo de monitoreo calcula señales no supervisadas y supervisadas. Las
señales no supervisadas comparan la partición de referencia con la ventana
actual usando solo $X$. En este grupo se incluyen PSI, KS, divergencia de
Jensen--Shannon, distancia de Hellinger y pruebas $\chi^2$ para variables
categóricas. Adicionalmente, Evidently se usa como herramienta de monitoreo
\textit{batch}; su \text{DataDriftPreset} permite evaluar cambios de
distribución entre dos conjuntos de datos y reportar drift por columna
\cite{evidently_data_drift,evidently_github}.

Las señales supervisadas se calculan cuando existen etiquetas disponibles. En
este grupo se ubican $\Delta AUC_t$, $\Delta \ell_t$, F1 y el flujo de errores
del modelo. River ADWIN se usa como detector en línea sobre el flujo de error de
predicción; ADWIN mantiene una ventana adaptativa y detecta cambios cuando las
medias de subventanas recientes dejan de ser compatibles \cite{river_adwin}. El
recorrido fila por fila de los datos tabulares se implementa con utilidades de
\text{river.stream} \cite{river_iter_pandas}.

El resultado del monitoreo se resume en un vector de observación para el agente:
promedios y máximos de métricas de drift por variable, indicadores de drift por
familia de variables, señales supervisadas de desempeño y costos históricos de
acciones previas. Estas señales no se interpretan como observaciones perfectas
del estado real; se consideran evidencia parcial del tipo, severidad y efecto
del drift.

% ------------------------------------------------------------
\section{Entorno de decisión y acciones del agente}
% ------------------------------------------------------------

El problema de actualización selectiva se modela como un POMDP, porque el agente
no observa directamente $P_t(X)$ ni $P_t(Y \mid X)$, sino métricas estimadas
sobre ventanas finitas \cite{kaelbling1998pomdp}. Cada episodio contiene entre
10 y 20 pasos de decisión. En cada paso se aplica un drift sintético, se calcula
el vector de observación, el agente selecciona una acción, se ejecuta la acción
y se mide si el rendimiento mejora frente al estado sin mantenimiento.

\begin{table}[H]
    \centering
    \caption{Acciones de mantenimiento consideradas por DARL}
    \label{tab:acciones_darl}
    \begin{tabular}{lll}
        \hline
        \textbf{Acción}          & \textbf{Componente afectado} & \textbf{Interpretación}     \\
        \hline
        \textit{Defer}           & Ninguno                      & Mantener el pipeline actual \\
        \textit{Update features} & $f_{\phi}$                   & Reajustar preprocesamiento  \\
        \textit{Update model}    & $g_{\theta}$                 & Reajustar modelo predictivo \\
        \textit{Retrain all}     & $f_{\phi}$ y $g_{\theta}$    & Reentrenar todo el pipeline \\
        \hline
    \end{tabular}
\end{table}

La hipótesis operacional es que distintos tipos de drift favorecen distintas
acciones. Un \textit{covariate shift} puro debería favorecer la actualización de
\textit{features} cuando el modelo sigue siendo válido. Un \textit{concept
    drift} debería favorecer la actualización del modelo. Un drift combinado puede
requerir reentrenamiento completo. La acción \textit{Defer} permite capturar
escenarios donde el cambio detectado no justifica el costo de mantenimiento.

% ------------------------------------------------------------
\section{Recompensa y entrenamiento PPO}
% ------------------------------------------------------------

La recompensa combina recuperación de rendimiento y costo computacional. Para
cada acción $a_t$, se compara el AUC obtenido después de aplicar la acción con
el AUC del mismo estado drifted sin mantenimiento. El costo se normaliza contra
el costo de reentrenar todo el \textit{pipeline}, usado como referencia de mayor
intervención:

\begin{equation}
    r_t(a_t)
    =
    \lambda_A
    \left(
    AUC_t^{(a_t)} - AUC_t^{(\text{drift})}
    \right)
    -
    \lambda_C
    \left(
    \frac{1}{2}
    \frac{T_t^{(a_t)}}{T_t^{(\text{all})}}
    +
    \frac{1}{2}
    \frac{M_t^{(a_t)}}{M_t^{(\text{all})}}
    \right).
    \label{eq:darl_reward}
\end{equation}

\noindent\textbf{Donde:}
\begin{itemize}
    \item $r_t(a_t)$ representa la recompensa recibida tras ejecutar la acción $a_t$ en el paso $t$.
    \item $AUC_t^{(a_t)}$ representa el AUC-ROC obtenido después de aplicar la acción seleccionada.
    \item $AUC_t^{(\text{drift})}$ representa el AUC-ROC del pipeline drifted sin mantenimiento.
    \item $T_t^{(a_t)}$ representa el tiempo de ejecución de la acción seleccionada.
    \item $T_t^{(\text{all})}$ representa el tiempo requerido por el reentrenamiento completo.
    \item $M_t^{(a_t)}$ representa el uso máximo de memoria de la acción seleccionada.
    \item $M_t^{(\text{all})}$ representa el uso máximo de memoria del reentrenamiento completo.
    \item $\lambda_A$ y $\lambda_C$ ponderan la importancia relativa de recuperar rendimiento y reducir costo.
\end{itemize}

Conceptualmente, esta recompensa evita que el agente elija siempre
\textit{Retrain all}. Una acción recibe mayor recompensa si recupera AUC con bajo
costo relativo; recibe menor recompensa si recupera poco rendimiento o si incurre
en un costo similar al reentrenamiento completo sin una mejora proporcional. Este
diseño sigue la motivación de mantenimiento consciente del costo: la mejor acción
no es necesariamente la que maximiza rendimiento absoluto, sino la que ofrece el
mejor balance entre recuperación y recursos \cite{mahadevan2024cost}.

El agente se entrena con PPO, descrito en la Sección~\ref{sec:ppo}. En esta
metodología, PPO no reemplaza al detector de drift; aprende una política sobre
las señales producidas por el monitor. Su función es mejorar gradualmente la
probabilidad de seleccionar acciones que producen mayor recompensa acumulada,
manteniendo actualizaciones de política estables mediante el objetivo recortado
\cite{schulman2017ppo}.

% ------------------------------------------------------------
\section{Métricas de evaluación}
% ------------------------------------------------------------

La evaluación final separa métricas de rendimiento y métricas de cómputo. Las
métricas de rendimiento incluyen AUC-ROC como métrica principal, F1-score como
medida complementaria y $\Delta AUC$ como caída o recuperación frente al estado
pre-drift. Las métricas de cómputo incluyen tiempo de actualización, memoria
máxima y costo relativo frente a \textit{Retrain all}. Esta separación permite
responder dos preguntas distintas: si la acción corrige el deterioro y si lo
hace con menos recursos que el reentrenamiento completo.

